% Generated draft: requires mathematical and source review.
\documentclass[11pt]{article}
\usepackage[margin=25mm]{geometry}
\usepackage{fontspec}
\setmainfont{Latin Modern Roman}
\usepackage{amsmath,amssymb,mathtools,unicode-math}
\setmathfont{Latin Modern Math}
\usepackage{enumitem,needspace,xurl,hyperref}
\setlength{\parindent}{0pt}
\setlength{\parskip}{5pt}
\setlength{\emergencystretch}{3em}
\newcommand{\N}{\mathbb{N}}
\newcommand{\Z}{\mathbb{Z}}
\newcommand{\Q}{\mathbb{Q}}
\newcommand{\R}{\mathbb{R}}
\newcommand{\C}{\mathbb{C}}
\newcommand{\F}{\mathbb{F}}
\newcommand{\A}{\mathbb{A}}
\newcommand{\PP}{\mathbb{P}}
\newcommand{\E}{\mathbb{E}}
\newcommand{\eps}{\varepsilon}
\newcommand{\dd}{\,\mathrm{d}}
\begin{document}
\section*{UnsolvedMath ID: 20002865}
\textbf{Source problem code:} catalogue entry
\subsection{Definitions and mathematical statement}
Let $\mathbb B^k$ be the unit ball in $\C^k$. X. Huang's source question takes a discrete subgroup $\Gamma\subset\operatorname{Aut}(\mathbb B^N)$ and a proper holomorphic embedding $f:\mathbb B^n\to\mathbb B^N$ with $\gamma(f(\mathbb B^n))\subset f(\mathbb B^n)$ for every $\gamma\in\Gamma$ and compact image quotient. It asks whether $f$ is linear. The natural rigidity interpretation is linear up to ball automorphisms, equivalently totally geodesic image:
\[
 \exists\alpha\in\operatorname{Aut}(\mathbb B^n),\ \beta\in\operatorname{Aut}(\mathbb B^N):
 \quad\beta\circ f\circ\alpha(z)=(z,0).
\]
Literal coordinate linearity without normalization is different, since even a ball automorphism need not be linear. The intended interpretation and lower bound on $n$ are not explicit.

Put $S=f(\mathbb B^n)$. Since $\Gamma$ is a group, applying the inclusion to $\gamma^{-1}$ gives $\gamma S=S$. The aid defines the induced homomorphism
\[
 \rho:\Gamma\to\operatorname{Aut}(\mathbb B^n),\qquad
 \rho(\gamma)=f^{-1}\gamma f,
\]
and reports a finite kernel, discrete image $\Lambda$, and compact $\mathbb B^n/\Lambda$. The finite kernel fixes an image point and lies in a compact ambient stabilizer; discreteness is essential to its finiteness. The image is thus a cocompact source lattice, not necessarily a cocompact lattice of the ambient ball. After a finite-index torsion-free passage in the ambient linear group, the kernel is removed and the equivariant embedding descends to a compact complex-hyperbolic source manifold inside an ambient ball quotient. The precise finite-index choice is part of the reduction, not an initial torsion-free assumption.

The aid applies this to report automorphism-equivalent linearity in the following cases: $N=n$; $1<n<N\le2n-1$ via Cao--Mok; algebraic image via Chan--Mok; a $C^2$ boundary extension; and, for $n\ge2$, a boundary H\"older extension with exponent greater than $1/2$ via a 2026 theorem. The full boundary-extension and algebraicity hypotheses are not supplied, so their precise application must be checked. In particular the dimension restrictions of the middle-codimension and H\"older assertions are retained.

In the intended higher-dimensional setting the aid leaves the nonalgebraic, boundary-irregular high-codimension regime $N\ge2n$ unresolved in its checked literature. The $n=1$ scope is not certified by the stated higher-dimensional theorems. The finite-kernel/lattice reduction supplies hypotheses for those rigidity results; it does not alone prove total geodesy in every dimension and codimension.
\subsection{Short English statement}
The group action is actually two-sided and induces a cocompact source lattice with finite kernel. The aid removes torsion and applies several qualified rigidity theorems. The general irregular high-codimension case remains beyond that reduction, and linearity requires an automorphism-normalized interpretation.
\subsection{Sources}
\begin{itemize}
\item Local UnsolvedMath registry, permanent ID 20002865, source code catalogue entry; created and updated 14 August 2026. Canonical \texttt{aim-several-complex-variables-notes.json}, zero-based \texttt{notes[22]}, tag \texttt{problem}, item 4 attributed to X. Huang.
\item AIM ``Emerging applications of complexity for CR mappings,'' 2010 workshop according to the supplied verification: \url{https://aimath.org/WWN/crmappings/crmappings.pdf}. Wording is marked exact; original remarks and literature field are empty.
\item Research attempt 2: \texttt{partial\_result}, problem status at run \texttt{partially\_solved}, rigor label \texttt{proved}, novelty confidence low. It names Cao--Mok, Chan--Mok, and an unidentified 2026 H\"older theorem without full titles, theorem statements, authors for the latter, or authoritative links.
\item Ulam AI / UnsolvedMath Contributors, CC BY 4.0. Machine-generated research aid, not peer reviewed; independently verify before citation.
\end{itemize}
\end{document}
